\clearpage
\section{The Appointed Texts}
Latin follows the 1962 Roman Missal, pp.~411--412. The Introit antiphon and the three orations use the historical English of the Cummiskey hand missal (1861), Nineteenth Sunday after Pentecost. Biblical English is Douay--Rheims--Challoner. It retains the Bible's wording where the Latin has a liturgical introduction or adaptation. These are historical study texts.

\subsection{Introit}\label{proper-introit}
\properlocus{Printed 411, n. 1679; Ps 77(78):1 after the antiphon.}
\begin{latinproper}
Salus populi ego sum, dicit Dominus: de quacumque tribulatione clamaverint ad me, exaudiam eos: et ero illorum Dominus in perpetuum. Ps. 77, 1 Attendite, popule meus, legem meam: inclinate aurem vestram in verba oris mei. V. Gloria Patri.
\end{latinproper}
\begin{englishwitness}{Cummiskey 1861}
I am the Saviour of my people, saith the Lord: in whatever distress they call on me, I will hear them: and will be their Lord for ever.
\end{englishwitness}
\begin{englishwitness}{Psalm verse: Douay--Rheims--Challoner}
Attend, O my people, to my law: incline your ears to the words of my mouth.
\end{englishwitness}

\subsection{Collect}\label{proper-collect}
\properlocus{Printed 411, n. 1680.}
\begin{latinproper}
Omnipotens et misericors Deus, universa nobis adversantia propitiatus exclude: ut mente et corpore pariter expediti, quae tua sunt, liberis mentibus exsequamur. Per Dominum.
\end{latinproper}
\begin{englishwitness}{Cummiskey 1861}
O Almighty and merciful God, favourably defend us from all adversity: that being free both in soul and body, we may with security of mind perform thy service. Thro’.
\end{englishwitness}

\subsection{Epistle}\label{proper-epistle}
\properlocus{Printed 411, n. 1681; Eph 4:23--28.}
\begin{latinproper}
Fratres: Renovamini spiritu mentis vestrae, et induite novum hominem, qui secundum Deum creatus est in iustitia et sanctitate veritatis. Propter quod deponentes mendacium, loquimini veritatem unusquisque cum proximo suo: quoniam sumus invicem membra. Irascimini, et nolite peccare: sol non occidat super iracundiam vestram. Nolite locum dare diabolo: qui furabatur, iam non furetur; magis autem laboret, operando manibus suis, quod bonum est, ut habeat unde tribuat necessitatem patienti.
\end{latinproper}
\begin{englishwitness}{Douay--Rheims--Challoner}
And be renewed in the spirit of your mind: And put on the new man, who according to God is created in justice and holiness of truth. Wherefore, putting away lying, speak ye the truth, every man with his neighbour. For we are members one of another. Be angry: and sin not. Let not the sun go down upon your anger. Give not place to the devil. He that stole, let him now steal no more: but rather let him labour, working with his hands the thing which is good, that he may have something to give to him that suffereth need.
\end{englishwitness}

\subsection{Gradual}\label{proper-gradual}
\properlocus{Printed 411, n. 1682; Ps 140(141):2.}
\begin{latinproper}
Dirigatur oratio mea sicut incensum in conspectu tuo, Domine. V. Elevatio manuum mearum sacrificium vespertinum.
\end{latinproper}
\begin{englishwitness}{Douay--Rheims--Challoner}
Let my prayer be directed as incense in thy sight; the lifting up of my hands, as evening sacrifice.
\end{englishwitness}

\subsection{Alleluia}\label{proper-alleluia}
\properlocus{Printed 411, n. 1683; Ps 104(105):1.}
\begin{latinproper}
Alleluia, alleluia. V. Confitemini Domino, et invocate nomen eius: annuntiate inter gentes opera eius. Alleluia.
\end{latinproper}
\begin{englishwitness}{Douay--Rheims--Challoner}
Give glory to the Lord, and call upon his name: declare his deeds among the Gentiles.
\end{englishwitness}

\subsection{Gospel}\label{proper-gospel}
\properlocus{Printed 411--412, n. 1684; Matt 22:1--14.}
\begin{latinproper}
In illo tempore: Loquebatur Iesus principibus sacerdotum et pharisaeis in parabolis, dicens: Simile factum est regnum caelorum homini regi, qui fecit nuptias filio suo. Et misit servos suos vocare invitatos ad nuptias, et nolebant venire. Iterum misit alios servos, dicens: Dicite invitatis: Ecce prandium meum paravi, tauri mei et altilia occisa sunt, et omnia parata: venite ad nuptias. Illi autem neglexerunt: et abierunt, alius in villam suam, alius vero ad negotiationem suam: reliqui vero tenuerunt servos eius, et contumeliis affectos occiderunt. Rex autem cum audisset, iratus est: et missis exercitibus suis, perdidit homicidas illos, et civitatem illorum succendit. Tunc ait servis suis: Nuptiae quidem paratae sunt, sed qui invitati erant, non fuerunt digni. Ite ergo ad exitus viarum, et quoscumque inveneritis, vocate ad nuptias. Et egressi servi eius in vias, congregaverunt omnes, quos invenerunt, malos et bonos: et impletae sunt nuptiae discumbentium. Intravit autem rex, ut videret discumbentes, et vidit ibi hominem non vestitum veste nuptiali. Et ait illi: Amice, quomodo huc intrasti non habens vestem nuptialem? At ille obmutuit. Tunc dixit rex ministris: Ligatis manibus et pedibus eius, mittite eum in tenebras exteriores: ibi erit fletus et stridor dentium. Multi enim sunt vocati, pauci vero electi.
\end{latinproper}
\begin{englishwitness}{Douay--Rheims--Challoner}
And Jesus answering, spoke again in parables to them, saying: The kingdom of heaven is likened to a king who made a marriage for his son. And he sent his servants to call them that were invited to the marriage: and they would not come. Again he sent other servants, saying: Tell them that were invited, Behold, I have prepared my dinner: my beeves and fatlings are killed, and all things are ready. Come ye to the marriage. But they neglected and went their ways, one to his farm and another to his merchandise. And the rest laid hands on his servants and, having treated them contumeliously, put them to death. But when the king had heard of it, he was angry: and sending his armies, he destroyed those murderers and burnt their city. Then he saith to his servants: The marriage indeed is ready; but they that were invited were not worthy. Go ye therefore into the highways; and as many as you shall find, call to the marriage. And his servants going forth into the ways, gathered together all that they found, both bad and good: and the marriage was filled with guests. And the king went in to see the guests: and he saw there a man who had not on a wedding garment. And he saith to him: Friend, how camest thou in hither not having on a wedding garment? But he was silent. Then the king said to the waiters: Bind his hands and feet, and cast him into the exterior darkness. There shall be weeping and gnashing of teeth. For many are called, but few are chosen.
\end{englishwitness}
The Creed follows the Gospel.

\subsection{Offertory}\label{proper-offertory}
\properlocus{Printed 412, n. 1685; adapted Ps 137(138):7.}
\begin{latinproper}
Si ambulavero in medio tribulationis, vivificabis me, Domine: et super iram inimicorum meorum extendes manum tuam, et salvum me faciet dextera tua.
\end{latinproper}
\begin{englishwitness}{Douay--Rheims--Challoner}
If I shall walk in the midst of tribulation, thou wilt quicken me: and thou hast stretched forth thy hand against the wrath of my enemies: and thy right hand hath saved me.
\end{englishwitness}
The Latin adds \textit{Domine} and gives the verbs of extending and saving in the future. The historical Bible English has the perfect forms printed above.

\subsection{Secret}\label{proper-secret}
\properlocus{Printed 412, n. 1686.}
\begin{latinproper}
Haec munera, quaesumus, Domine, quae oculis tuae maiestatis offerimus, salutaria nobis esse concede. Per Dominum nostrum Iesum Christum, Filium tuum: Qui tecum vivit et regnat in unitate.
\end{latinproper}
\begin{englishwitness}{Cummiskey 1861}
Grant, we beseech thee, O Lord, that the offerings we bring before thy divine Majesty may avail to our salvation. Thro’.
\end{englishwitness}
The Preface of the Most Holy Trinity is appointed. The Latin conclusion above retains the longer form printed here; the two other orations retain their shorter printed abbreviations.

\subsection{Communion}\label{proper-communion}
\properlocus{Printed 412, n. 1687; Ps 118(119):4--5.}
\begin{latinproper}
Tu mandasti mandata tua custodiri nimis: utinam dirigantur viae meae, ad custodiendas iustificationes tuas.
\end{latinproper}
\begin{englishwitness}{Douay--Rheims--Challoner}
Thou hast commanded thy commandments to be kept most diligently. O! that my ways may be directed to keep thy justifications.
\end{englishwitness}

\subsection{Postcommunion}\label{proper-postcommunion}
\properlocus{Printed 412, n. 1688.}
\begin{latinproper}
Tua nos, Domine, medicinalis operatio, et a nostris perversitatibus clementer expediat, et tuis semper faciat inhaerere mandatis. Per Dominum.
\end{latinproper}
\begin{englishwitness}{Cummiskey 1861}
May the healing efficacy of these thy mysteries, O Lord, mercifully free us from our perverseness, and make us always obedient to thy commandments. Thro’.
\end{englishwitness}

